\chapter{Introduction}

\startcontents[chapters]
\printmyminitoc 



\epigraph{\itshape  Everyone knows what attention is. It is the taking possession by the mind, in clear and vivid form, of one out of what seems several simultaneously possible objects or trains of thought.}{-- William James}

Attention is a field widely discussed and studied in neuroscience, psychology, cognitive science and machine learning \citep{chun2011taxonomy,cho2015describing}. Attention is the process of selectively focusing on a discrete aspect of information while ignoring other perceivable information. A widely accepted feature of attention is that it facilitates efficient use of the available computational resources.

Over time, the understanding of attention has progressed significantly, leading to the development of various models aimed at elucidating its mechanisms. The hypothesis put forth by \citet{descartes2010passions}, suggesting that the pineal body controls attention, is now considered invalid, and alternative models have been developed to offer more accurate explanations. Early investigations focused on psychophysics, particularly exploring the range of visual attention.  For instance, \citet{Helmholtz1989} introduced the concept of covert attention, which proposes that visual attention can operate independently of eye movements. While initial approaches viewed attention as a top-down guided internal state,  Gestalt theory \citep{kohler1947gestalt} presented attention as a bottom-up computation, emphasizing the role of the focus of attention. These early models challenged Descartes' (1649) outdated hypothesis that the pineal body controls attention.

Neurophysiological findings by \citet{Sechenov1863} suggested that inhibition from the central nervous system plays a crucial role in attention control. This discovery spurred the development of computational models by \citet{tsotsos1995modeling} and theoretical frameworks by \citet{Pavlov1927} that incorporated inhibition as a core process in attention. According to Pavlov, unexpected stimuli capture attention (facilitation), while stimuli in nearby cortical areas are suppressed through inhibitory processes. Further exploration of the role of inhibition as an essential attention process was conducted by \citet{itti2009bayesian}. Notably, attenuation theory by \citet{treisman1980feature} proposed that unattended signals are attenuated rather than filtered, reconciling opposing viewpoints. These models and theories contributed to the understanding of attention as a mechanism for coping with limited information processing capacity \citep{broadbent1958perception}.

Broadbent's comprehensive attention model, grounded in cognitive psychology, viewed attention as a biological mechanism to cope with the limited capacity of information processing. This model sparked a debate regarding the stage at which selection occurs: early in the information pipeline, as proposed by Broadbent, or later in the selection process, as suggested by \citet{norman1968toward}. Treisman introduced the concept of attenuating unattended signals as a middle-ground perspective \citep{treisman1980feature}. \citet{milner1974model} proposed that attention not only selects relevant features but also provides feedback to early stages of information processing. This framework was incorporated into Adaptive Resonance Theory by \citet{grossberg1975neural}. Subsequent neurophysiological evidence demonstrated the top-down influence of the attentional state on the activation of perceptual circuitry, indicating that feedback can occur at any stage of the information pipeline.


Many computational models of visual saliency originate from Treisman et al. Feature Integration Theory (FIT) of spatial visual attention proposed in 1980. The FIT aimed to explain the performance difference between pop-out stimuli and conjunction search. Texton theory by \citet{bergen1983parallel} demonstrated that certain features allow for rapid discrimination of a target and surrounding outliers (pop-out), while others do not. According to the FIT, a saliency map, often referred to as a ``master map,'' integrates information from separate feature maps to identify salient locations. Treisman and Gelade also addressed the binding issue, which involves the cohesive representation of an object by binding its features (color, shape, location, etc.) together. \citet{koch1987shifts} proposed a computational implementation of the FIT, where the saliency map is a weighted sum.


Visual attention is a specific form of attention that operates within the visual modality. It involves the selective processing and allocation of attentional resources to visual stimuli. Visual attention allows us to prioritize and focus on specific visual features, objects, or regions within our visual field while suppressing or inhibiting the processing of other visual inputs. It can operate at different levels. At the early perceptual level, it involves the selection and processing of basic visual features such as color, shape, and motion. This early selection process helps filter out irrelevant visual information and enhance the salience of relevant visual stimuli. At a higher cognitive level, visual attention enables us to selectively attend to objects or regions of interest, guiding our gaze and directing our focus within the visual scene.

The mechanisms of visual attention include both bottom-up and top-down processes. Bottom-up attention is driven by salient or physically distinctive features of stimuli that automatically capture our attention, such as a bright color or sudden movement. Top-down attention, on the other hand, is driven by our goals, expectations, and prior knowledge. It allows us to voluntarily direct our attention to specific stimuli based on their relevance or importance in a given context. Visual attention plays a crucial role in various cognitive processes, such as visual perception, object recognition, scene understanding, and visual search tasks. It helps us efficiently process and interpret visual information, guiding our interactions with the visual world.

% attention, selectivity and its link to neuroscience
The cognitive science literature depicts several aspects of attention, such as it can be concentrated, it can focus on a particular modality, it can be divided, it can be selective, and it can have a finite capacity. However, selectivity is its most characteristic feature. Selective attention is necessary because of the limited availability of resources. Visual attention and selective attention are closely interconnected concepts that involve the cognitive process of allocating attentional resources to specific stimuli while ignoring or suppressing others. Visual attention refers to the selective processing and prioritization of visual information, while selective attention encompasses the broader ability to attend to stimuli across different sensory modalities.

Selective attention is the cognitive process of focusing on one or a limited number of sensory stimuli while disregarding irrelevant inputs. Various theories have been proposed to explain selective attention, including bottleneck theories and load theories. Bottleneck theories, such as Filter Theory \citep{broadbent1958perception}, Late Selection Theory \citep{deutsch1963attention}, and Attenuation Theory \citep{treisman1964selective}, focus on the flow and filtering of information. Load theories, like Perceptual Load Theory by \citet{lavie1994perceptual}, Dilution Theory \citep{tsal2010diluting} address the allocation of perceptual and cognitive resources. However, operationalizing these constructs and validating the theories can be challenging. Selective attention is essential for daily functioning, preventing overload of the information processing system. 

%% expanding selective attention
Early theories of attention, such as Donald Broadbent's Filter Theory, proposed a "bottleneck" model of selective attention, likened to a bottle with a narrow opening. According to Broadbent, stimuli enter a sensory buffer where their physical characteristics are assessed, allowing only a few to pass through the selective filter. Unselected stimuli decay in the buffer, while the selected ones proceed to be processed for meaning and determine how we respond. Broadbent used the dichotic listening task to study selective attention, finding that participants performed better when attending to one ear at a time. However, criticisms arose regarding where stimuli gain meaning within the attention process, with the cocktail party effect suggesting that analysis occurs before filtering.

Deutsch and Deutsch proposed the late selection theory as an alternative to address limitations of Broadbent's theory. They suggested that all stimuli are analyzed for meaning, but only selected ones pass the filter based on their physical characteristics and relevance. Anne Treisman introduced the attenuation theory, suggesting that stimuli are not filtered but attenuated or enter the sensory register at a lower intensity, gaining meaning early on. Her theory addresses the limitation of Broadbent's theory regarding the cocktail party effect. Treisman used the dichotic listening task with complete words and found that people often combine messages from both ears, implying that the unattended message still holds meaning regardless of retention. These early theories laid the groundwork for understanding selective attention and the processing of stimuli based on their relevance and physical characteristics.

% relate visual attention and selective attention
Visual attention involves the mechanisms and processes that enable us to focus on relevant visual stimuli and filter out irrelevant or distracting visual inputs. It allows us to direct our attention to specific regions or objects within the visual field, selectively process their features, and integrate them into our perceptual experience. Visual attention plays a crucial role in various tasks, such as visual search, where we actively scan our environment to find a specific target among distractors. Selective attention, on the other hand, extends beyond the visual domain and encompasses attentional processes across different sensory modalities, including auditory, tactile, and cognitive inputs. It involves the ability to prioritize and allocate attentional resources to relevant stimuli or information while disregarding or suppressing irrelevant or less important stimuli from all sensory channels.

While visual attention is a specific subset of selective attention that focuses on the processing and filtering of visual information, it is interconnected with other modalities of attention. For example, during a complex task that requires both visual and auditory processing, selective attention allows us to prioritize the relevant visual stimuli while simultaneously attending to relevant auditory cues or instructions.

%%%
Recently visual attention has gained tremendous attention in the field of artificial intelligence. 
Visual attention helps in answering \textit{what} to look and \textit{where} to look. It has been vastly studied in psychology and neuroscience \citep{posner1990attention,cohen1990control,phaf1990slam,bundesen1990theory,desimone1995neural,mozer1998computational,corbetta2002control,o2006making,petersen2012attention,moore2017neural} and more recently by \citet{flesch2022orthogonal,dekker2022curriculum}. These studies have acted as a source of inspiration for several artificial intelligence models \citep{khosla2007bio,lindsay2018biological} including the ones proposed in this thesis. 

% types of visual attention
There are three categories of selectivity in a visual attention system: by spatial location \textit{(space-based)} \citep{posner1980orienting,posner1982neural}, by object membership \textit{(object-based)}~\citep{duncan1984selective,egly1994shifting,vecera1994does,kramer1997object} and by particular features of the input \textit{(feature-based)}~\citep{harms1983color,driver1989movement,kramer1991perceptual,baylis1992visual,duncan1996objects}. 

\paragraph{Visual Spatial Attention}

% overt-covert attention,
Every second, our eye makes small and rapid movements several times, known as saccades. These eye movements change the locus of attention. Visible shifts of attention, such as saccades, are known as \textit{overt} visual attention. One more method used to emphasize a spatial location without any over-the-shift of the fovea location is \textit{covert} attention. An example is the subject's fixation on a particular region throughout a task where the stimulus is likely to appear. This region is also referred to as the ``\textit{spotlight}'' of attention. Certain visual patterns that involve edges, contrast, or motion automatically attract attention. These patterns are known as ``\textit{salient}'' \citep{itti2001computational}. In the presence of task-specific information, these saccadic movements are controlled in a top-down fashion around the particular visual target instead of the salient regions. Eye movements are one of the possible ways to control visual attention. 


% feature based attention
\paragraph{Visual Feature Attention} 
When the focus of attention is on features like color, shape or orientation instead of location, it is known as feature-based attention. It is an example of covert visual attention. Cueing the right features enhances the system's performance. It is used in tasks such as visual search combining covert feature-based attention with overt attention. Feature-based attention is global as opposed to spatial attention, i.e., when attention is focused on a particular feature, neurons representing that particular feature in the visual space are also modulated~\citep{saenz2002global}. It is related to object attention, i.e., instead of attending to an abstract feature, the attention is deployed at a specific object in a visual scene~\citep{chen2012object}. A single feedforward pass in the visual hierarchy can segregate the objects of a visual scene if there is a distinct salient difference between them as opposed to a complex scene where recurrent and serial processing might be required~\citep{lamme2000distinct}. 

% bottom-up vs. top down attention
In addition to feature-based or spatial attention, another widely accepted classification is characterized by the type of data processing~\citep{connor2004visual,buschman2007top}. There are two types of data processing, \textit{bottom-up} and \textit{top-down}. In a bottom-up attention process, external factors guide the attentional process because of their inherent properties, like their color or sudden motion in the scene. It is fast and primitive sensory driven. In top-down attention, there is an internal attentional guidance mechanism based on prior knowledge and current goals, like searching for food if one is hungry. It can ignore the salient stimuli and focus on the target object or event. 

% attention, working memory and the central controller
Attention is also involved while performing tasks requiring multiple sensory signals. In the presence of multiple tasks or sensory signals, the central executive controller helps to route the focus of attention. The Central executive controller is responsible for coordinating activity with the cognitive system for directing attention, decision making and maintaining task goals. Context and history are deemed helpful to executing tasks optimally -- making it highly related to the working memory. Attention is furthermore seen as the output of the central controller. The controller selects the targets of attention and passes them to the system responsible for its implementation. There is a three-way relationship between executive control, working memory and attention in such a way that the focus of attention is selected by the executive controller based on the contents of the working memory ~\citep{soto2008automatic}. Although all the objects in the working memory can influence attention, the executive controller helps decide which one should affect the most~\citep{olivers2011difference}. These vast and extensive cognitive studies related to attention have inspired the field of AI and helped to boost its performance (Figure~\ref{fig:atnml}).


\begin{figure}[htbp]
\centering
  \includegraphics[width=.7\linewidth]{images/intro/atttention_ml.png}
  \caption{A summary of attention in Cognitive science and machine learning (\href{https://www.frontiersin.org/articles/10.3389/fncom.2020.00029/full\#h5}{source})} \label{fig:atnml}
\end{figure}

% history of attention in AI (before 2000)
The first attempt to adapt the attention mechanism in a neural network was made in the 1980s when the improved version of \textit{Neocognitron} \citep{Fukushima1980} incorporated selective attention \citep{fukushima1987neural} to decompose the image into elementary features. Later, \citet{Fukushima1993} modified the network to recognize and segment characters in cursive handwriting. \citet{Postma1997} proposed an attentional scanning model, \textit{SCAN}, to attend to and identify object patterns without decomposing the scene into elementary features. As an alternative to these static neural approaches, \citet{Schmidhuber1991LearningTG} proposed a sequential model inspired by the sequential eye movements for object detection. In this model, a neural controller learns sequential generation of fovea trajectories to reach the target. Furthermore, data processing types inspired the development around the same time, thereby leading to a model extracting the region of interest using bottom-up and top-down processing~\citet{MilaneseIntegration1994}.

% history of attention in AI (during 2000s)
By the early 2000s, the influence of attention on the evolution of neural networks increased. \citet{miau2001neural} proposed a model of primate vision integrating both, \textit{what} and \textit{where} pathways. The model has a fast visual attention-based front-end to select the most salient image areas and a slow back-end to recognize objects in those selected areas. Another model based on the primate selective mechanism is presented in \citet{salah2002selective} with the idea of selectively attending to relevant parts of the input image. In this model, a neural network analyzes the input image and generates posterior probabilities for the Markov models. Attention has also been used for object recognition \citep{walther2002attentional} and scene analysis \citep{schill2001scene}. 

% history of attention in AI (during 2015) attention in NLP: neural machine translation
The year 2015 marks the new beginnings of attention-based architectures with the introduction of the attentional model for Neural Machine Translation (NMT) \citep{Bahdanau2014-vr,luong-etal-2015-effective} and image captioning \citep{xu2015show}. In NMT, the expectation is to learn continuous representations of variable-length sequences. \textit{Recurrent neural networks} (RNNs) like LSTMs \citep{10.1162/neco.1997.9.8.1735}, GRUs \citep{cho-etal-2014-properties} and Quasi-RNNs \citep{DBLP:conf/iclr/0002MXS17} were some of the popular sequence models for representation learning at that time. While these RNNs' output depends on the previous elements in a sequence, traditional feedforward neural networks assume that inputs and outputs are independent of each other. Nonetheless, their limitation includes their inability to parallelize computations -- making them slow during training and their fixed-size memory -- bottleneck for long-range interactions \citep{vaswani2017attention}. 

% history of attention in AI (during 2015) RNN based models, description
Models used for NMT typically consist of encoder-decoder architecture \citep{cho-etal-2014-learning}. Typically, both encoder and decoder are RNN, where the encoder takes an input sequence of fixed-length vector and represents it again to another fixed-length vector. A decoder then takes this encoded vector to generate the output sequence token by token. However, this method has two challenges; first, the encoder compresses the input sequence into a fixed vector length which may lead to the loss of information \citep{cho-etal-2014-properties}. Second, the model is incapable of aligning between input and output sequences which is essential for tasks such as translation or summarization \citep{Young2018trends}. While generating the output sequence, the decoder also lacked the mechanism to selectively focus on relevant input tokens. Later, \citet{bahdanau2014neural} proposed a sequence-to-sequence modeling task with the help of soft attention, emphasizing the parts of the sentence relevant to predicting the target word. \citet{bahdanau2014neural} extended the basic encoder-decoder by letting the model search a set of input words while generating target words. It allowed the model to focus on information needed to generate the subsequent target sequence. 

% history of attention in AI (during 2015) RNN based models, description
In the following two years, the adoption of attentional mechanisms in neural networks diversified. Content-based soft attention mechanism \citep{goodfellow2016deep} is used in Neural Turing Machine (NTM) \citep{graves2014neural} with end-to-end training. Around the same time, \citet{cheng2016long} used a form of attention called intra-attention in the Long Short-Term Memory (LSTM) \citep{hochreiter1997long} architecture. \citep{hochreiter1997long} embedded a memory network inside the LSTM architecture to store the contextual representation of the input. This memory network has a set of key and value vectors in the hidden state to represent what is stored in the memory. These vectors are used to estimate the intra-attention with the previously stored tokens in the memory as opposed to the self-attention mechanism used by \citet{vaswani2017attention} where interaction between the whole input sequence is estimated. One of the first uses of the self-attention mechanism in NLP is done by \cite{parikh-etal-2016-decomposable}. 

Since then, self-attention mechanisms have become an integral part of sequence modeling allowing the network to model dependencies between input and output sequences irrespective of their distances. A self-attention layer calculates a single-shot interaction between all pairs of words in a sequence. 

\section{Self-attention-based Transformer architecture}
% history of attention in AI (during 2017) transformer based models - detailed description
In 2017 \citet{vaswani2017attention} proposed a novel architecture, \textit{Transformer} for NLP. It is predominantly a self-attention network driving the waves of advances in AI. A Transformer architecture (Figure~\ref{fig:vaswani}) includes a stack of encoder and decoder blocks. Each encoder block is identical and contains a self-attention layer and a feedforward layer. The encoder's input flows through the self-attention layer helping the encoder to look at other words while encoding the current word. Its output is then fed to the feedforward layer. The same feedforward network is applied independently to each word. While a decoder consists of an encoder-decoder attention block in addition to the self-attention layer and a feedforward layer helping the decoder to focus on relevant parts of the input sequence.

\begin{figure}[ht]
\centering
  \includegraphics[width=1\linewidth]{images/intro/transformer.png}
\caption{Transformer architecture proposed by \citet{vaswani2017attention} (\href{https://deepfrench.gitlab.io/deep-learning-project/}{source})} 
\label{fig:vaswani}
\end{figure}

In the NLP task, each word of the input sequence is first converted into an embedding vector. They are provided as input to the encoder block, passing through a self-attention layer and feedforward network. The obtained output vector is fed to the next encoder block. Using the self-attention layer, the Transformer models the relationship between the current word with other relevant words of a sequence. 

% history of attention in AI (during 2017) self-attention in transformers
In a self-attention layer, its input vector is transformed into a key ($K$), query ($Q$) and value ($V$) vectors of dimension $d_q$ = $d_k$ = $d_v$ = 512 using a learnable matrix transformation. At first, the score ($S$) is calculated to determine the amount of focus to place on the other words in a sequence while encoding the current word. This score is calculated using the dot product between the query and key vectors ($S~=~Q~.~K^T$). It is normalized ($S=S/\sqrt{d_k}$) to stabilize the gradients, and later, using a $softmax$, converted into probabilities. The extent of the probability score shows the relevancy of the current word with other words in the sequence. This score is multiplied by the value vector ($V$) so that relevant words are given additional focus while irrelevant words are neglected in the subsequent layers. 

\[ Attention~(Q,K,V)~ = softmax(\frac{Q~.~K^T}{\sqrt{d_k}})~.~V\]

The self-attention mechanism proposed by \cite{vaswani2017attention} has an additional feature called \textit{multi-head} attention (MHA). It helps to improve the performance in two ways: by augmenting the network's ability to focus on multiple positions and by giving distinct representational subspaces to each word. For example, if there are eight heads, eight sets of K, Q and V matrices exist, each representing a unique representational subspace. They are concatenated before passing through the feedforward network (Figure~\ref{fig:intromha}). 

\begin{figure}[htbp]
\centering
  \includegraphics[width=1\linewidth]{images/intro/transformer_multi-headed_self-attention-recap.png}
\caption{Illustration of Multi-head attention mechanism in a Transformer network (\href{https://jalammar.github.io/images/t/transformer_multi-headed_self-attention-recap.png}{source})} 
\label{fig:intromha}
\end{figure}

The key characteristic of NLP tasks is the order of the words in a sequence. However, the operations we have discussed till now are permutation invariant. Positional embedding vectors are added to each input embedding vector to address this issue. These vectors help the model estimate the position of each word in a sequence in the projection space (i.e., $K/Q/V$).

Positional encoding in the Transformer is an active and vibrant research area. Vanilla Transformer \citep{vaswani2017attention} uses absolute positional encoding; however, more recent work \citep{DBLP:journals/corr/abs-1810-04805,Dosovitskiy2020-iq} prefers a learned \citep{gehring2017convolutional} or relative positional encoding \citep{shaw2018self}. The absolute coordinate system does not encode translational equivariance, while relative geometry could. \cite{ramachandran2019stand,bello2019attention} studied different positional encoding techniques and established that relative positional encoding offers the best results while providing additional advantages like encoding for an unseen length of sequences (refer to \cite{wu2021rethinking} for a review). An overview of different positional encoding strategies used in NLP is discussed by \citet{dufter2021position}. 


The residual connection around the self-attention layer and a feedforward network is an essential module in an encoder block. It is followed by the layer normalization step~\citep{baevski2018adaptive,wang2019learning,Dosovitskiy2020-iq}. A residual connection is added to each sub-layer in the encoder (and decoder) to strengthen the flow of information and achieve higher performance.

At the end of the encoding steps, decoding begins. The decoder uses the key ($K$) and value ($V$) vectors from the top-most encoder block for its encoder-decoder attention layer. It helps to focus on the appropriate locations of the input sequence. At each step, the decoder layer provides an element of the final output sequence. This output vector is again fed to the subsequent decoder layer in the next time step. This process continues till the end of the sequence. An independent set of positional encoding is applied on the decoder side. 

The encoder-decoder attention module is similar to the multi-head self-attention mechanism described earlier. The only difference is that the key $K$ and value $V$ vectors are obtained from the top-most encoder block, and the query vector $Q$ is derived from the previous self-attention layer of the decoder. Unlike the encoder, self-attention layers in the decoder are only allowed to access previously obtained output by masking the future words of a sequence. Masking future positions is done to prevent the decoder from cheating during the training phase -- otherwise, it will already know what is coming next. The linear layer at the end of the decoder block is a fully connected neural network. It projects the vector obtained from the decoder layers into a logit vector. This logit vector represents the complete vocabulary of the language where translation has to be performed. A softmax converts this logit to the probabilities and represents the concerned word from the available vocabulary. 

In terms of computational complexity, for a sequence of length $n$ and dimensionality $d$, self-attention layers are faster than recursive or convolutional layers when $n$ is smaller than $d$, which is typically the case.

\begin{table}[ht]
\begin{center}
\vspace{.1in}
\begin{tabular}{cccc}
\hline
\textbf{Layer Type}     & \textbf{Complexity} & \textbf{Sequential } & \textbf{Maximum} \\
        & \textbf{per Layer} & \textbf{operations} & \textbf{ path length} \\\hline
Self-Attention & $\mathcal{O}$($n^2.d$)  & $\mathcal{O}$(1)    & $\mathcal{O}$(1)      \\
Recurrent      & $\mathcal{O}$($n.d^2$)  & $\mathcal{O}$(n)    & $\mathcal{O}$(n)      \\
Convolutional  & $\mathcal{O}$($k.n.d^2$) & $\mathcal{O}$(1)   & $\mathcal{O}$($log_k$(n))  \\\hline                      
\end{tabular}
\caption{Complexity comparison of different networks for a sequence of length $n$, kernal size $k$ and dimensionality $d$ \citep{vaswani2017attention} }
\label{tab:complexity}

\end{center}
\end{table}

\section{Self attention in vision tasks}

\paragraph{Why self-attention for vision?}

In a self-attention mechanism, each word of a sequence is correlated with all the others. Thus containing information about the rest of the sequence -- increasing the receptive field size equivalent to the length of a sequence. In some sense, images are no different from NLP sequences. Computer vision can take inspiration from the NLP domain to model long-range interactions between pixels with the added benefit of multi-head attention helping to parallelize these interactions. With the help of the multi-head attention method, different heads can focus on modeling different relations between pixels. For example, in a visual reasoning task where the objective is to count the number of pairs of shapes in an image, one head can focus on finding a pair while the other can focus on counting them. It helps the network to model self-similarity within an image. Images such as natural scenes and paintings display a great amount of self-similarity. Such non-local self similarity property was earlier explored for applications such as texture synthesis \citep{efros1999texture}, object detection and segmentation \citep{Wang_2018_CVPR}, bilateral filtering \citep{tomasi1998bilateral} and image classification \citep{pmlr-v80-parmar18a}. Hereafter, the main focus of this thesis will be computer vision.  

\paragraph{Self-attention with CNN}

In a computer vision task, the resolution of the images could reach around 1000$\times$1000 px. Applying a self-attention mechanism to all these pixels ($~10^6$ in number) is computationally expensive because of the quadratic complexity associated with the length of the sequence. Convolutional layers, on the other hand, do not have this bottleneck. However, they face trouble capturing long-range interactions because of their inability to scale up with the large receptive fields. 

To address this problem, there are predominantly two approaches. The first is to reduce the self-attention operation cost to a linear scale. Aligned to this line of work, \citet{ramachandran2019stand} proposed a pure stand-alone attention model for vision tasks by replacing the convolution operations with self-attention operations. Nonetheless, the self-attention operation used in this approach is local. Another similar linear attention variant Halo~\citep{vaswani2021scaling} uses block-wise local attention to improve speed and accuracy. 

The second approach is to build hybrid CNN-Transformer architectures where the convolutions operations are used to encode the input image, and attention is applied to those encoded features. \citet{srinivas2021bottleneck} explored a hybrid combination of CNNs and multi-head self-attention (MHSA) models and showed that replacing the $3\times 3$ kernal size convolutional layer in the bottleneck blocks of ResNet \citep{he2016deep} with MHSA layers improved several CNN baselines. Interestingly, DETR \citep{carion2020end} showed that concatenating the Transformer model at the end of the feature-extraction network is helpful for tasks like detection, localization, and segmentation. 


There are four broad categories of research to incorporate self-attention mechanism with CNN, which are as follows:

\textbf{Inserting few attention modules in between residual blocks}: Along this line of work, \citet{Wang_2018_CVPR,chen20182} proposed a non-local block similar to \citet{ramachandran2019stand} and used them for video-based applications. In this network, features are gathered and propagated motivated by the squeeze and excite \citep{hu2018squeeze} network. As mentioned earlier, these methods only focus on the spatial dimension for calculating the non-local interaction, so \citet{yue2018compact} added a correlation factor between the channels to improve the model effectiveness. Similarly, \citet{Shen_2021_WACV} proposed a method to bring down the quadratic complexity of the self-attention mechanism to a linear scale. We demonstrate a unique way to incorporate self-attention with a feedforward network in \hyperref[chapter1b]{Chapter 3} where the intermediate features of the network are passed through the self-attention layer to find the global association. This attention is applied directly over the feature space in contrast to the previously used methods of squeezing the feature vector dimensions to save computations. 

\textbf{Inserting attention modules at the end}: Usually, such models have a front-end of convolutional block acting as a feature extraction module for self-attention block as back-end. These models are used for tasks like object detection and semantic segmentation. \citet{huang2019ccnet} designed criss-cross attention that learns the complete image dependency recurrently for semantic segmentation tasks using dot-product attention. Moving away from this trend of using self-attention operations, \citet{carion2020end} proposed DETR architecture by placing a Transformer model as the back-end. 

\textbf{Replacing convolution layers by self-attention layers}: Self-attention mechanism used in this line of research is primarily local in nature to decrease the computational demand associated with the increasing sequence length in an image which is directly proportional to total pixel count. \citet{bello2019attention} made a unique attempt to augment the feature maps of convolutional layers with the self-attention modules. Feature maps obtained with the help of the self-attention module are concatenated with the feature maps of CNNs. They discovered that replacing all the feature maps of CNN with the feature maps of self-attention layers degrades the system's performance. Contrary to their finding, \citet{ramachandran2019stand} came up with the architecture replacing all the convolution layers with a local self-attention layer and achieved better performance than a fully convolutional network on the image classification task. 

In addition to these four categories of research, where the primary focus is on computer vision applications, cognitive studies also explore self-attention mechanisms. In one of the first studies by \citet{whittington2022relating}, neural representations of Hippocampal formation are related to the Transformer model. They did this correspondence with the help of the Tolman-Eichenbaum Machine \cite{whittington2020tolman}, a model for hippocampal formation. This work showed that when recurrent positional encodings are used in Transformer, they replicate spatial representations of hippocampal formations like place cells and grid cells. To analyze from an attentional point of view, we studied the role of a self-attention layer of the Transformer model in understanding visual reasoning tasks in \hyperref[chapter1b]{Chapter 3}. This layer is used as a feature-based or spatial attention layer. A multi-head self-attention layer is significantly different from the other existing self-attention models where the span of attention in the dot-product mechanisms is local. They proposed a self-attention mechanism that could be applied globally over a feedforward network's spatial or feature space. This method gives the network higher representation power because of its ability to use multi-head attention. We also built a cognitive architecture inspired by the active vision literature relating to the shifting of the spotlight of attention in \hyperref[chapter2b]{Chapter 4}. This attention routing is implemented with a controller module consisting of an attention module and an LSTM layer, which generates a query to guide the shifting. More studies in the NLP domain focus on relating language models to brain activations; however, a similar trend is yet to be seen in the computer vision domain. 

These developments exploring the self-attention mechanisms propelled toward building a fully self-based attention architecture for computer vision applications. Evolutions in the NLP domain were vital in inspiring the fully self-attention-based architecture for vision tasks. 


\section{Transformer-based vision architecture}


The first fully self-attention-based Transformer architecture is presented by \citet{Dosovitskiy2020-iq}. It is known as Vision Transformer (ViT) (Figure~\ref{fig:vit}). In this architecture, an input image is divided into a sequence of image patches called visual tokens and transforms those patches before passing them to the network. The core idea is to treat each pixel as a token and pass it to the Transformer network. However, with the increasing size of the number of pixels, attention cost scales quadratically, so patches of 16$\times$16 pixels are used instead. Each patch is flattened and linearly projected to a vector of the desired dimension. As the network is agnostic to the positions of these patches w.r.t. the input image, position embeddings are added to learn the 2D structure. ViT learns this encoded structural information while training. A learnable class embedding token is also added at the beginning of the sequence. A class embedding token is inspired from \citet{devlin2018bert} that is learned along with other patches while training the network. This learnable token eventually helps to predict the classification label with the help of a multi-layer perceptron (MLP) head. 


\begin{figure}[htbp]
\centering
  \includegraphics[width=1\linewidth]{images/intro/vit.png}
\caption{Vision Transformer architecture \citep{Dosovitskiy2020-iq}} 
\label{fig:vit}
\end{figure}

When the ViT is trained on a mid-sized dataset like ImageNet \citep{deng2009imagenet}, outcomes are not impressive because of their lack of inductive biases such as translational equivariance and locality. ViT experiences difficulty learning image-specific inductive biases like a CNN, as the model never sees the complete 2D image during training but only a sequence of transformed patches. Such CNN-like biases are compensated by training the model with massive databases like JFT-300M and fine-tuning it for downstream tasks. ViT learns the spatial relationship from scratch, which raises its demand for extra training data and longer training time. \cite{touvron2021training} proposed DeIT, equalizing this pre-training-related bottleneck using techniques like the teacher-student distillation approach and robust augmentation methods. DeIT, when trained on ImageNet by incorporating these methods, surpasses the performance of the ViT model. 

Vision transformers are a front-runner in capturing the long-range dependencies in an image, yet they fail to account for local features as CNNs do. A wide gap is perceived between ViT and CNN learnability. \cite{wu2021cvt,guo2022cmt,yuan2021incorporating,graham2021levit,dai2021coatnet,Peng_2021_ICCV} analysed the potential weaknesses in directly applying Transformer model from NLP domain and proposed a combination with convolutional network. \cite{wu2021cvt} proposed a Convolutional vision Transformer (CvT) and presented a convolutional-based patch projection of image tokens along with a hierarchical design. Another alternative, LocalViT \citep{li2021localvit} proposed depthwise convolution to capture local features.
Meanwhile, LeViT \cite{graham2021levit} enhanced the inference speed of ViT by designing multistage transformer architecture and downsampling the image using attention. In yet another network proposed by \cite{zhou2021elsa}, it incorporated locality without convolutions with the help of enhanced local self-attention using Hadamard attention and ghost head. Hadamard attention is more computational-friendly than dot-product attention, while ghost heads increase the channel capacity by combining attention maps. 

A striking network, ConViT, proposed by \cite{d2021convit} took a step further to incorporate the convolutional biases into the Transformer architecture. \cite{d2021convit} initialized self-attention layers with soft convolutions with the help of Gated-Positional-Self-Attention (GPSA). This self-attention block is characterized by locality strength and head-specific center of attention. The locality factor determines how much the head should focus around its center of attention. For any given query patch, which head should give attention to which position is decided by the head-specific center of attention. With suitable parameters setting, ConViT can have ViT-like expressive power and could be trained in low-data regimes like CNNs. In a collaborative project with paleobotanist, we proposed \textit{conviformer} \citet{vaishnav2022conviformers} to incorporate convolutional biases into any vision transformer with minimal architectural change. With the conviformer architecture, the network can also attend to higher-resolution images and provide compatibility with the base architecture used. 

\paragraph{Challenges}

\begin{figure}[htb]
\centering
  \includegraphics[width=.7\linewidth]{images/intro/Transformer-computational.png}
  \caption{Computational demands for training Transformers vs. CNNs. Compute needed to train a Transformer network has increased by 275 times in the last two years. (\href{https://blogs.nvidia.com/blog/2022/03/25/what-is-a-transformer-model/}{source})} \label{fig:transcomp}
\end{figure}

Transformer architecture confronts a two-front challenge. It requires enormous data for training to learn the right inductive bias and the computational cost associated with the sequence length. Figure~\ref{fig:transcomp} compares the computational requirement of different Transformers and CNNs models. An empirical study is done by \citet{zhai2022scaling} on the scalability of the ViT. They report that scaling up training samples and parameters of the model scale up the overall performance of the model; nevertheless, this plateaus quickly for smaller models as they cannot leverage additional data. It indicates that larger models have the scope to improve their representation learning abilities. Training a Transformer model requires massive data to compete with inductive biases like translation invariance similar to CNN. The self-attention mechanisms in a Transformer learn such image-specific concepts during longer training times, thereby significantly increasing the compute requirements. Strong data augmentation techniques nowadays compensate for the vast dataset requirement. 
% Visualizing and interpreting Transformer models is also a challenging task.


Transformer architecture furthermore lacks explicit mechanisms to attend to local neighborhoods. A commonly accepted solution to this issue is to restrict the attention mechanism to the local area \cite{parmar2018image} or to incorporate structural priors on attention like Sparsity \citep{child2019generating}. It makes a dense attention matrix into a sparse matrix limiting the computations. Regardless, the approach has some limitations. Sparse matrix multiplication operations are uncommon for hardware accelerators. 

An additional computational bottleneck is calculating the dot product operation in the self-attention layer. Existing techniques to handle this situation are half-precision, gradient accumulation and gradient checkpointing. Tensor computations on modern hardware architectures are effectively done with 16-bit float tensors. Sometimes higher precision is required while calculating the loss, which doubles the required memory. This precision handling is carried out with the help of \textit{apex} library\footnote{\url{https://nvidia.github.io/apex/}}. On a fixed GPU/TPU machine, a large model may only fit a single-digit batch size, ultimately leading to unstable learning. A multivariate chain rule is used to incorporate the dynamics related to bigger batch sizes. It sums the gradients for a larger batch and computes the gradient descent at the end. For more bigger models, the trade-off is to separate the model into different chunks and compute the gradient in a forward/backward pass for each chunk.   

Our proposed \textit{conviformer}  \cite{vaishnav2022conviformers} addressed ViT's inability to process longer sequences which restricted ViT to smaller resolution images. In the conviformer, the input image is passed through a convolutional backbone, down-sampling the image to 224$\times$224 (a commonly accepted input resolution). With the help of a convolutional frontend, the network makes sure to introduce the inductive biases of CNN into the network. The feature vectors obtained by the CNN modules are later passed to the base architecture of the vision transformer. This technique holds the compatibility of the network with the base model and provides a performance boost with insignificant additional computational cost.

Finally, training a huge Transformer model has negatively impacted the environment. Compute cost and the complexity associated with the Transformer are directly related to environmental factors such as $CO_2$ emission \citep{Strubell2020} and high energy consumption \citep{You2020Drawing}. There is also a cost associated with mining rare metals for manufacturing these hardware accelerators. 

\section{Original Contributions}
Our contributions are as follows:
\begin{itemize}
    \item We present a novel fine-grained taxonomy for the SVRT tasks by systematically analyzing the ability of feedforward neural networks.
    \item We first propose a self-attention-augmented feedforward network modeled as spatial or feature-based attention. 
    \item Our attentional networks analysis on SVRT tasks provides a granular computational account of visual reasoning and yields testable neuroscience predictions regarding the differential need for feature-based versus spatial attention depending on the type of visual reasoning problem.
    \item Next, we present a novel end-to-end trainable guided-attention module to learn to solve visual reasoning challenges in a data-efficient manner.
    \item We show that our guided-attention module learns to shift attention to task-relevant locations and gate relevant visual elements into a memory bank; 
    \item We show that our architecture demonstrate zero-shot generalization ability and learns compositionally. GAMR is capable of learning efficiently by re-arranging previously-learned elementary operations stored within a reasoning module. 
    \item Our architecture sets new benchmarks on two visual reasoning challenges, SVRT~\citep{fleuret2011comparing} and ART~\citep{Webb2021EmergentST}. 
\end{itemize}

The work presented in \hyperref[chapter1a]{Chapter 2} and \hyperref[chapter1b]{Chapter 3} are taken from our following publication:
\begin{itemize}
    \item \textbf{Mohit Vaishnav}, Remi Cadene, Andrea Alamia, Drew Linsley, Rufin VanRullen, Thomas Serre; ``Understanding the Computational Demands Underlying Visual Reasoning.'' \textit{Neural Computation} 2022; 34 (5): 1075–1099. doi: \url{https://doi.org/10.1162/neco_a_01485}
\end{itemize}

The work presented in \hyperref[chapter2b]{Chapter 4} is taken from our following publication:
\begin{itemize}
    \item \textbf{Mohit Vaishnav}, Thomas Serre. ``GAMR: A Guided Attention model for (visual) Reasoning.'' \textit{International Conference on Learning Representations (ICLR)} 2023, \href{https://openreview.net/forum?id=iLMgk2IGNyv}{https://openreview.net/forum?id=iLMgk2IGNyv}
\end{itemize}